% Transcription — fonds Alexandre Grothendieck, shelfmark 62, batch 4 (pages 61-80).
% Established from the facsimile archives/batches/62/batch-04.pdf (a copy of
% 62.pdf fetched through the project relay, no cover sheet: PDF page k =
% archivists' page 60+k), per the skill transcribe-grothendieck. The folder is
% ninety-six pages in five batches; this is the fourth.
%
% Pass: Opus 5.5 (claude-opus-5-5), 2026-09-26 — first pass, unchecked against the pages by a human.
%
% Pages skipped: none. All twenty pages carry his mathematics (tan paper,
% ink; no typed versos in this batch). Pages 61-73 are in a fair, regular
% hand; pages 74-80 are on greyer paper in a faster hand.
%
% Pages summarised: none in full. Page 76 is cancelled by long diagonal
% strokes and read only in fragments; it is given by fragments, with one
% \note for the cancellation. Page 73 is half blank.
%
% His pagination: none; the only page numbers are the archivists' pencilled
% ones bottom left. No dates. Numbered runs of his own, all restarting:
% p. 61 « 1. fonction \wp z » and a second heading « 1.(?) fonction \zeta z »;
% pp. 65-69 the Legendre sn run, formulas (1)-(13) (1°, 2°, 3° on pp. 65-66);
% pp. 70-73 the same material for sn_1, cn_1, dn_1, formulas (1)-(9);
% p. 72 a third recapitulation of sn, formulas (1)-(6); pp. 74-80, under the
% heading « Fonctions elliptiques » (p. 74), sections 1) généralités with
% items a')-g), 2) « Les relations entre \wp, \wp' » (p. 78), 3) « Fonction
% \zeta(u) » (p. 80), with formula runs (1)-(3) on p. 74, (3) on p. 77,
% (1)-(3) on p. 78, two boxed formulas on p. 79 whose numbers read badly
% ((4)/(7), (5)/(8)), then (6), (7) and (7)-(11) on p. 80. Recorded in notes
% where they fall.
%
% What the batch is about. Classical transcendental theory of elliptic
% functions, continuing the brown-ink run of pp. 59-60 (batch 3, under the
% cover « Théorie transcendante (Généralités classiques) », Théorèmes I-V).
% Pages 61-64: Weierstrass \wp (definition, parity, Laurent expansion, the
% differential equation \wp'^2 = 4\wp^3 - g_2\wp - g_3, e_1, e_2, e_3,
% \wp'', \wp'''), \zeta (quasi-periods \eta_i, Legendre's relation), Hermite's
% decomposition, \sigma as a Weierstrass product and its quasi-periodicity,
% Jacobi's formula (poles and zeros), addition formulas. Pages 65-73:
% Legendre's sn z as inverse of the integral of dZ/\sqrt{(1-Z^2)(1-k^2Z^2)},
% expressed by \wp and by \sigma, periods 2\omega_1, 4\omega_2, 4\omega_3,
% poles, zeros and residues; cn, dn, derivation formulas and quadratic
% relations; then two re-draftings (sn_1 with k_1^2 = (e_3-e_1)(e_2-e_1),
% p. 70-71, 73; a recapitulation of sn, p. 72). Pages 74-80: a fresh,
% more conceptual draft « Fonctions elliptiques »: the field E of elliptic
% functions for the period lattice, identified with meromorphic functions on
% C/\mathfrak{P}; \wp by Mittag-Leffler; even elliptic functions are
% rational in \wp, polynomial when poles lie in the lattice; [E : C<\wp>] = 2
% and f = R_1(\wp) + \wp' R_2(\wp); the coefficients c_k as Eisenstein sums;
% zeros \omega_i of \wp', distinct e_i; \wp'^2 = 4\wp^3 - g_2\wp - g_3 =
% 4\prod(\wp - e_i); c_3, c_4 in terms of c_1, c_2; \zeta and \eta.
%
% Notation fixed once for the batch: his period lattice, a large gothic
% capital, as \mathfrak{P}; his C<\wp> (angle brackets) as
% \mathbf{C}\langle\wp\rangle; primed sums as \mathop{{\sum}'}; the
% involution f -> f(-z) as \check{f}, as he writes it. Inside formulas an
% unreadable index or number is written « ? », as in batches 1-3.
%
% Hardest pages: 75-77 (the fast connective prose of 1) c)-g)), 76 (the
% cancelled page), 79 (the prose around the two boxed equations).
% Readings to check first: the denominator of sn \omega_2 on p. 67 (given
% as « ? »); the signs in (5) of p. 72 (sn(z+2\omega_2) = sn z as written);
% the first coefficient of (8) on p. 80 (c_2/3 as written, c_1/3 expected);
% the formula numbers on p. 79.

\documentclass[11pt,a4paper]{article}
% Le préambule commun à toutes les transcriptions du fonds Grothendieck.
%
% Il tient en une idée : **l'incertitude se compose, elle ne se résout pas.**
% Une transcription qui lisse un mot douteux a détruit la seule chose pour
% laquelle on la faisait — un lecteur doit pouvoir savoir, à chaque endroit,
% ce qui était sur le feuillet et ce qui a été ajouté. Les six macros qui
% suivent sont donc l'appareil critique, et non de la décoration.
%
% Les mêmes commandes sont reconnues par `scripts/render.mjs`, qui produit la
% vue de lecture du site. Une macro ajoutée ici doit l'être aussi là-bas,
% faute de quoi le rendu échouera bruyamment — ce qui est voulu : un rendu
% silencieusement faux est pire qu'un rendu absent.

\ProvidesPackage{grothendieck}[2026/08/08 Transcriptions du fonds Grothendieck]

\usepackage{amsmath,amssymb,amsthm}
\usepackage{mathrsfs}
\usepackage{stmaryrd}
% Les diagrammes commutatifs sont partout dans ce fonds : c'est la langue
% même de la géométrie algébrique de Grothendieck, et une transcription qui
% les rendrait en prose ne transcrirait rien.
\usepackage{tikz-cd}
% Français par défaut — c'est la langue des feuillets — mais l'anglais est
% chargé aussi : la traduction et le résumé sont des documents anglais, et la
% typographie française (espaces avant les deux-points) y serait une faute.
% Ces documents basculent par \selectlanguage{english} en tête de corps.
\usepackage[english,french]{babel}
\usepackage{fontspec}
\usepackage{geometry}
\geometry{a4paper,margin=25mm,marginparwidth=28mm}
\usepackage{marginnote}
\usepackage{xcolor}
% ulem, not soul. `soul`'s \st and \ul are built on a character-by-character
% scan that XeLaTeX's font machinery breaks: the document fails with an
% undefined control sequence rather than degrading. ulem does the same two jobs
% and is Unicode-engine safe. `normalem` stops it hijacking \emph, which the
% transcriptions use for Grothendieck's own emphasis.
\usepackage[normalem]{ulem}
\usepackage{hyperref}
\hypersetup{colorlinks=true,linkcolor=black,urlcolor=black,pdfborder={0 0 0}}

% --- Métadonnées du lot -----------------------------------------------------
% Reprises telles quelles par le script de rendu, qui les affiche en tête de
% la vue de lecture. Elles fixent ce que le fichier prétend couvrir : sans
% elles, une transcription flotte sans point d'attache dans un fonds de seize
% mille feuillets.

\newcommand{\folder}[1]{\gdef\grothFolder{#1}}
\newcommand{\batch}[1]{\gdef\grothBatch{#1}}
\newcommand{\leaves}[2]{\gdef\grothFirst{#1}\gdef\grothLast{#2}}
\newcommand{\foldertitle}[1]{\gdef\grothFoldertitle{#1}}
\gdef\grothFolder{?}\gdef\grothBatch{?}\gdef\grothFirst{?}\gdef\grothLast{?}\gdef\grothFoldertitle{}

% --- Le résumé --------------------------------------------------------------
%
% Il ouvre la lecture modernisée, sous le titre « Résumé », et il est composé
% en petit. Ce n'est pas un ornement : c'est le seul passage du document qui
% ne soit pas des mathématiques, et un lecteur doit pouvoir le reconnaître
% avant de l'avoir lu — puis le sauter s'il connaît déjà le sujet, ou s'y
% arrêter s'il ne le connaît pas. Le filet qui le suit marque la reprise du
% corps.

\newenvironment{resume}
  {\small\setlength{\parskip}{0.45em}}
  {\par\medskip\hrule\medskip}

% --- Mots-clés ---------------------------------------------------------------
%
% En anglais, en clôture du résumé de la lecture modernisée : le vocabulaire
% moderne sous lequel chercher ce que les feuillets construisent. Le manifeste
% du site (`npm run manifest`) les extrait pour étiqueter la cote entière —
% cette ligne est la source unique des tags, il n'y a pas de fichier à côté.
\newcommand{\keywords}[1]{\par\smallskip\noindent
  {\footnotesize\textsc{Keywords} --- \textit{#1}}\par}

% --- L'appareil critique ----------------------------------------------------

% Le numéro de feuillet, en marge. C'est l'ancre qui permet de régler un
% désaccord sur une formule en regardant *un* feuillet plutôt qu'une chemise
% de six cents. Le site s'en sert aussi pour tourner les pages du fac-similé
% quand on fait défiler la transcription.
\newcommand{\leaf}[1]{%
  \marginnote{\footnotesize\color{gray!70}\textsf{#1}}%
  \label{leaf:#1}\ignorespaces}

% Illisible : on ne devine pas. Un mot inventé qui se lit comme les autres est
% le pire résultat possible.
\newcommand{\ill}{\textcolor{red!60!black}{[\,\ldots\,]}}

% Lecture douteuse : le mot est donné, et signalé comme incertain.
\newcommand{\uncertain}[1]{\uline{#1}}

% Ajout de l'éditeur — une lettre manquante, un mot sous-entendu.
\newcommand{\add}[1]{\textcolor{blue!50!black}{[#1]}}

% Biffé par Grothendieck lui-même. Ce qu'il a rayé fait partie du document :
% une rature dit souvent où la pensée a changé de direction.
\newcommand{\struck}[1]{\sout{#1}}

% Note du transcripteur, sur l'état matériel du feuillet.
\newcommand{\note}[1]{\par\smallskip{\small\color{gray!60!black}\textit{[#1]}}\par\smallskip}

% Note marginale de Grothendieck, distincte du corps du texte.
\newcommand{\marginal}[1]{\par\smallskip\noindent\hspace{1em}%
  {\small\color{gray!60!black}#1}\par\smallskip}

% --- Titre ------------------------------------------------------------------

\AtBeginDocument{%
  \begin{center}
    {\large\bfseries Fonds Alexandre Grothendieck --- cote n\textsuperscript{o}~\grothFolder}\\[2pt]
    {\itshape\grothFoldertitle}\\[4pt]
    {\small feuillets \grothFirst--\grothLast\ (lot \grothBatch)}\\[2pt]
    {\footnotesize\color{gray!60!black}Transcription établie d'après le fac-similé numérisé
     par l'Université de Montpellier.\\
     Les lectures incertaines sont soulignées, les illisibles notées [\,\ldots\,],
     les ajouts entre crochets.}
  \end{center}
  \vspace{1em}}

\endinput


\folder{62}
\batch{4}
\pages{61}{80}
\dating{[à partir de 1963]}
\watermark{Édition de démonstration}
\foldertitle{Courbes elliptiques (généralités - vieille rédaction) : notes manuscrites (s.d.)}

\begin{document}

\page{61}
\emph{1. fonction $\wp z$}
\note{titre souligné ; le numéro « 1. » est de sa main}

Définition intrinsèque : c'est une fonction elliptique d'ordre $2$, admettant
l'origine pour pôle double, \struck{y ayant pour \ill{}} \add{la fonction}
$\wp z - \frac{1}{z^2}$ s'annulant à l'origine. Si $\omega_1$ et $\omega_2$ sont
donnés, cela détermine \uncertain{sans plus} cette fonction elliptique.

Elle est paire [car $\varphi(z)-\varphi(-z)$ est elliptique, n'a pas de pôle,
s'annule à l'origine] \struck{et vérifie l'équation} \add{et est donc de la forme}
\[
\wp z = \frac{1}{z^2} + c_1 z^2 + c_2 z^4 + \cdots
\]
\note{à gauche de la formule, un début de formule biffé ; les indices et les
exposants de $c_1z^2 + c_2z^4$ sont petits et se lisent mal}
Elle vérifie l'équation
\[
\wp'^2 z = 4\wp^3 z - g_2\wp z - g_3 \qquad g_2 = 20c_1 \quad g_3 = 28 c_2
\]
C'est donc la fonction inverse de
\[
z = \int_\infty^{Z} \frac{dZ}{\sqrt{4Z^3 - g_2 Z - g_3}}
\]
On en déduit que si on pose
\[
\wp\omega_1 = e_1 \qquad \wp\omega_2 = e_2 \qquad \wp\omega_3 = e_3
\]
$e_1, e_2, e_3$ sont les racines de l'équation $4Z^3 - g_2Z - g_3 = 0$ ; on a
\uncertain{donc}
\[
\wp'\omega_1 = \wp'\omega_2 = \wp'\omega_3 = 0 .
\]
\struck{$\wp'\omega_1 = e_1$}\note{formule biffée, lue sous réserve}

On a
\[
\wp'' z = 6\wp^2 z - \frac{g_2}{2} \qquad \wp''' z = 12\,\wp z\,\wp' z
\]
\note{la seconde formule est soulignée en pointillé}
$\wp^{(2n)}z$ s'exprime en fonction de $\wp z$ par un polynôme de degré $n+1$ ;
$\wp^{(2n+1)}z$ par le produit par $\wp' z$ d'un polynôme en $\wp z$ de degré $n$.
\add{Développement en série double}
\note{écrit sous la ligne et appelé par un signe d'insertion après
« $\wp^{(2n+1)}z$ par »}

\emph{\uncertain{1.} fonction $\zeta z$.}
\note{titre souligné ; le chiffre ressemble au « 1. » du premier titre}
S'introduit par l'intégration des fonctions elliptiques. La fonction
$\int_0^z \bigl(\wp z - \frac{1}{z^2}\bigr)dz$ est uniforme, en lui ajoutant
$\frac{1}{z}$ on obtient une primitive de $\wp z$) ; on pose
\[
\zeta z = \frac{1}{z} - \int_0^z \Bigl(\wp z - \frac{1}{z^2}\Bigr)dz
\qquad \text{on a} \qquad \zeta' z = -\wp z
\]
\note{ainsi sur la page : « primitive de $\wp z$ », suivi d'une parenthèse
fermante sans ouvrante ; la formule qui suit donne $\zeta' = -\wp$}

\page{62}
$\zeta z$ est une fonction impaire, \struck{\ill{} périodique, \ill{}. On a}
\add{admettant les $w$ pour pôles simples (avec $1$ comme résidu)}.
\note{l'ajout est écrit au-dessus de la ligne ; « avec $1$ comme résidu » est
repris sous la ligne, dans la marge gauche}
\[
\zeta(z+2\omega_1) - \zeta z = \text{c}^{\text{te}} = 2\zeta\omega_1 = 2\eta_1
\]
on définit de même $\eta_2$ et $\eta_3$.

On établit par une intégration de \uncertain{$\int \zeta z\,dz$} sur un contour
convenable
\[
\omega_2\eta_1 - \omega_1\eta_2 = \frac{\pi i}{2} \quad \text{etc.} ;
\qquad \text{d'ailleurs} \quad \eta_1 + \eta_2 + \eta_3 = 0 .
\]

La fonction $\zeta z$ permet de calculer l'intégrale
\[
\int \frac{Z\,dZ}{\sqrt{4Z^3 - g_2Z - g_3}}
\]
on pose $Z = \wp z$ d'où
\[
\Bigl|\ \int \wp z\,dz = -\zeta z + \text{c}^{\text{te}}
\]
\note{le trait vertical renvoie à l'intégrale qui précède}

\struck{\ill{} \ill{} \ill{},}

\emph{Formule de décomposition de Hermite.} Soit $f(z)$ une fonction elliptique
de périodes $2\omega_1$ et $2\omega_2$ ; $a, b, \ldots$ ses pôles, distincts à une
période près deux à deux,
\[
\frac{A_\alpha}{(z-a)^\alpha} + \cdots + \frac{A_1}{z-a} \ ; \qquad
\frac{B_\beta}{(z-b)^\beta} + \cdots + \frac{B_1}{z-b} \ ; \qquad \ldots
\]
la partie principale de ses pôles ; on a alors
\[
\begin{aligned}
f(z) = \text{c}^{\text{te}} &+ \Bigl\lbrace A_1\zeta(z-a) + A_2\wp(z-a)
- \frac{A_3}{2!}\wp'(z-a) + \cdots
+ (-1)^\alpha \frac{A_\alpha}{(\alpha-1)!}\wp^{(\alpha-2)}(z-a) \Bigr\rbrace \\
&+ \Bigl\lbrace B_1\zeta(z-b) + \cdots \Bigr\rbrace + \cdots
\end{aligned}
\]
\note{suit une ligne de tirets qui sépare la section suivante}

\emph{Fonction $\sigma z$}\note{titre encadré et souligné}
Nous saurons intégrer toute la fonction elliptique au moyen de fonctions connues,
si nous savons intégrer $\zeta z\,dz$, son intégrale est \uncertain{complexe},
\struck{et} admet pour les point $w$ pour points logarithmiques. Il suffira de
chercher $\sigma z$ tel que
\[
\int \zeta z\,dz = \log \sigma z
\]
$\sigma z$ s'annulant aux points $w$.

\page{63}
On peut montrer directement que
\[
\sigma z = e^{\int \zeta z\,dz}
\]
est uniforme, \struck{est donc} \add{et même} une fonction entière,
\struck{\uncertain{il est}} on peut aussi procéder directement :

On a \struck{en effet}
\[
\wp z = \frac{1}{z^2} + \sum\mathop{{\sum}'} \left[\frac{1}{(z-w)^2} - \frac{1}{w^2}\right]
\]
la \uncertain{série étant} uniformément convergente pour $z$ borné, à distance
$>\varepsilon$ de $w$. Les \uncertain{dérivées} successives s'obtiennent par
dérivation terme à terme, la convergence est comme précédemment.

Notons que l'on obtient ainsi les coefficients du développement de $\wp z$
\[
c_1 = 3\sum\mathop{{\sum}'} \frac{1}{w^4} \qquad
c_2 = 5\sum\mathop{{\sum}'} \frac{1}{w^6} \qquad \cdots
\]
\note{l'exposant de $w$ dans $c_1$ est petit et peu net}
On a
\[
\zeta z = \frac{1}{z} + \sum\mathop{{\sum}'} \left[\frac{1}{z-w} + \frac{1}{w}
+ \frac{z}{w^2}\right]
\]
\note{le dénominateur du dernier terme est surchargé (un $2w^2$ corrigé en
$w^2$, semble-t-il)}
Posons enfin
\[
\sigma z = z \prod\mathop{{\prod}'} \Bigl(1 - \frac{z}{w}\Bigr)
e^{\frac{z}{w} + \frac{z^2}{2w^2}}
\]
Le produit \struck{est} converge uniformément \add{et absolument} si $z$ est borné ;
$\sigma z$ est une fonction entière, qui s'annule pour les $z = w$ et pas d'autres
points. Elle a été construite de sorte à avoir
\[
\frac{\sigma' z}{\sigma z} = \zeta z
\]
(À remarquer que, si on voulait construire a priori une fonction entière
s'annulant aux points $w$, aussi simple que possible, c'est $\sigma z$ qu'on
aurait trouvé)

$\sigma z$ est une fonction paire.\note{ainsi sur la page ; $\sigma$ est en fait
impaire} On a
\[
\sigma(z+2\omega_1) = -e^{2\eta_1(z+\omega_1)}\sigma z \qquad
\sigma(z+2\omega_2) = -e^{2\eta_2(z+\omega_2)}\sigma z
\]
de même pour $\omega_3$.

\page{64}
\emph{Formule de Jacobi.} Soient $a, b, \ldots, l$ les pôles
\struck{d'une} \add{d'une} fonction elliptique de pér.\ $2\omega_1$ et $2\omega_2$
(distincts à une pér.\ près) $a_1, b_1, \ldots, l_1$ les zéros ; \ldots
\[
a + b + \cdots + l = a_1 + b_1 + \cdots + l_1 + \Omega .
\]
$\Omega$ est une \supplied{période}.\note{la phrase s'arrête au bord droit sur
« $\Omega$ est une »}
On a alors
\[
f(z) = \text{c}^{\text{te}} \times
\frac{\sigma(z-a_1)\sigma(z-b_1)\cdots\sigma(z-l_1-\Omega)}
{\sigma(z-a)\sigma(z-b)\cdots\sigma(z-l)}
\]
soit
\[
f(z) = \text{c}^{\text{te}} \times
\frac{\sigma(z-a_1)\sigma(z-b_1)\cdots\sigma(z-l_1)}
{\sigma(z-a)\sigma(z-b)\cdots\sigma(z-l+\Omega)}
\]
soit encore, \ill{}
\[
f(z) = \text{c}^{\text{te}} \times e^{2(n_1\eta_1+n_2\eta_2)z}\,
\frac{\sigma(z-a_1)\cdots\sigma(z-l_1)}{\sigma(z-a)\cdots\sigma(z-l)}
\]
en posant $\Omega = 2n_1\omega_1 + 2n_2\omega_2$.
\note{plusieurs lettres des deux premières fractions sont surchargées}

\emph{Formules d'addition}
\[
\wp u - \wp v = -\frac{\sigma(u+v)\,\sigma(u-v)}{\sigma^2 u\,\sigma^2 v}
\]
\[
\zeta(u+v) = \zeta u + \zeta v + \frac{1}{2}\,\frac{\wp' u - \wp' v}{\wp u - \wp v}
\]
\[
\wp(u+v) + \wp u + \wp v = \frac{1}{4}\left(\frac{\wp' u - \wp' v}{\wp u - \wp v}\right)^2
= \bigl[\zeta(u+v) - \zeta u - \zeta v\bigr]^2
\]
\note{le dénominateur du coefficient $\frac14$ est lu $4$ sous réserve}

\page{65}
\emph{Fonction de Legendre $\mathrm{sn}\,z$}\note{titre souligné}

La fonction
\[
\int \frac{dZ}{\sqrt{(1-Z^2)(1-k^2Z^2)}}
\]
est tentant, car elle paraît très symétrique, et rappelle des intégrales
élémentaires liées aux fonctions trigonométriques et hyperboliques ($k=0$, $k=1$).

Appelons $\mathrm{sn}\,z$ l'intégrale de l'équation
\[
\Bigl(\frac{dZ}{dz}\Bigr)^2 = (1-Z^2)(1-k^2Z^2)
\]
qui prend la valeur zéro au point $z=0$, avec $Z'(0) = +1$. L'unicité de cette
solution se montre de diverses façons.

\emph{1°} En faisant un changement de variable convenable, on obtient
\[
Z = \frac{12\wp(z-z_0) - (5-k^2)}{12\wp(z-z_0) - (5k^2-1)}
\qquad
\begin{array}{l}
e_1 = \frac{1}{6}(1+k^2)\\[2pt]
e_2 = -\frac{1}{12}(k^2+6k+1)\\[2pt]
e_3 = -\frac{1}{12}(k^2-6k+1)
\end{array}
\]
\note{au dénominateur de cette première formule, le $\wp$ semble porter un
accent ; il ne le porte pas dans la formule (1) qui suit}
L'intégrale particulière est
\[
(1)\quad \mathrm{sn}\,z = \frac{12\wp(z-z_0) - (5-k^2)}{12\wp(z-z_0) - (5k^2-1)}
\qquad \text{avec} \quad \wp z_0 = \frac{5-k^2}{12} \quad
\wp' z_0 = -\frac{1-k^2}{2} .
\]
Ces formules sont de forme trop lourde, eu égard à d'autres plus simples. Mais on
voit l'existence de $\mathrm{sn}\,z$, qui est une fonction elliptique dont les
périodes s'expriment par des intégrales.

\page{66}
\emph{2°} On obtient
\[
Z^2 = \frac{1}{k^2}\bigl[\wp(z-z_0) - e_1\bigr]
\qquad (2)\ \left|\
\begin{array}{l}
e_1 = -\frac{1+k^2}{3}\\[2pt]
e_3 = \frac{2k^2-1}{3}\\[2pt]
e_2 = \frac{2-k^2}{3}
\end{array}\right.
\]
\note{les dénominateurs $3$ de $e_1$ et de $e_3$ sont mal formés (le premier
écrit sur un autre chiffre, le second proche d'un $7$)}
La solution \uncertain{envisagée} est
\[
(3)\quad \boxed{\ \mathrm{sn}^2 z = \frac{1}{k^2}\bigl[\wp(z-\omega_1) - e_1\bigr]\ }
\]

\emph{3°} Avec la même fonction elliptique, on a une solution. \uncertain{Finale}
\[
Z^2 = \frac{1}{\wp(z-z_0) - e_1}
\]
et la solution particulière est
\[
(4)\quad \boxed{\ \mathrm{sn}^2 z = \frac{1}{\wp z - e_1}\ }
\]
On montrerait facilement l'uniformité de \struck{$Z(z)$} $\mathrm{sn}\,z$ (les
racines de $\wp z - e_1 = 0$ étant double)

Les deux formules impliquent une identité :
\[
(\wp z - e_1)\bigl(\wp(z-\omega_1) - e_1\bigr) = k^2
\]
facile à vérifier, en passant à la fonction $\sigma$, on en \struck{montrant} que
le premier membre qui est elliptique, n'a pas de pôle, et est donc constant --
ceci est vrai pour \struck{tte} fonction $\wp z$ -- on trouve d'ailleurs
incidemment la formule
\[
\frac{e^{2\eta_1\omega_1}}{\sigma^4\omega_1} = 3e_1^2 - \frac{g_2}{4}
\]
\struck{, qui permet le calcul de $\eta_1$ et de \ill{}}
\struck{puis $\eta_2$ et $\eta_3$}

Dans le cas actuel, on a
\[
k^2 = \frac{e^{2\eta_1\omega_1}}{\sigma^4\omega_1} = 3e_1^2 - \frac{g_2}{4}
\]
\note{à droite, quelques croix ; l'exposant de $\sigma$ est lu $4$ sous réserve}
Toujours est-il que nous avons $\mathrm{sn}^2(z+\omega_1)\,\mathrm{sn}^2 z =
\frac{1}{k^2}$ ; si nous choisissons $\omega_1$ et \uncertain{$k$}
convenablement nous aurons
\[
(5)\quad \boxed{\ \mathrm{sn}(z+\omega_1)\,\mathrm{sn}\,z = \frac{1}{k}\ }
\]
\note{dans $\frac{1}{k^2}$, le numérateur est surchargé}

\page{67}
Expression au \emph{moyen de la fonction $\sigma$}
\[
(6)\quad \boxed{\ \mathrm{sn}\,z = \sigma\omega_1\, e^{\eta_1 z}\,
\frac{\sigma z}{\sigma(z+\omega_1)}
= -\sigma\omega_1\, e^{-\eta_1 z}\,\frac{\sigma z}{\sigma(z-\omega_1)}\ }
\]
\struck{$\mathrm{sn}\,z$ admet \ldots\ $\sigma(z$}\note{au-dessus du « $\mathrm{sn}$ »
biffé, un « $q$ » isolé}

On en déduit
\[
(7)\quad
\left\lbrace
\begin{array}{l}
\mathrm{sn}(z+2\omega_1) = \mathrm{sn}\,z\\
\mathrm{sn}(z+2\omega_2) = -\mathrm{sn}\,z\\
\mathrm{sn}(z+2\omega_3) = -\mathrm{sn}\,z\\[4pt]
\text{d'où}\quad \mathrm{sn}(z+4\omega_2) = \mathrm{sn}\,z\\
\phantom{\text{d'où}}\quad \mathrm{sn}(z+4\omega_3) = \mathrm{sn}\,z
\end{array}
\right.
\]
\note{dans la deuxième ligne, le « $\mathrm{sn}$ » est écrit sur un autre mot}
Donc $\mathrm{sn}\,z$ admet les périodes $2\omega_1$ et $4\omega_2$ (et
$4\omega_3$).

Les pôles sont $z = \omega_1 + 2n_1\omega_1 + 4n_2\omega_2$, ses zéros
$z = 2n_1\omega_1 + 4n_2\omega_2$; \note{ainsi sur la page : les indices des pôles
et des zéros se lisent mal ; on attendrait des pôles en $\omega_2$}
pôles et zéros sont simples ; $\mathrm{sn}\,z$ est une fonction d'ordre $2$. Les
zéros sont homologues de $0$ ou $2\omega_1$, où la dérivée est respectivement $1$ et
$-1$ ; les pôles sont homologues des pôles $\omega_2$ et $\omega_2+2\omega_1$, dont
les résidus sont resp.\ $\frac{1}{k}$ et $-\frac{1}{k}$ (en supposant toujours
choisi \add{le signe de} $\frac{1}{k}$ tel que l'on ait (5), ou, ce qui revient au
même
\[
k = -\frac{e^{\eta_1\omega_2}}{\sigma'\omega_2}\ \Bigr)
\]
\note{ainsi sur la page : $e^{\eta_1\omega_2}$ et $\sigma'\omega_2$, indices lus
sous réserve}

On voit sur 3 ou 4 que $\mathrm{sn}^2 z$ est pair ou impair ; elle n'est pas
\struck{im}paire, car on aurait $\mathrm{sn}\,\omega_2 = \mathrm{sn}\,\omega_2$,
mais aussi $\mathrm{sn}\,\omega_1 = -\mathrm{sn}\,{-\omega_1}$, d'où
$\mathrm{sn}\,\omega_1 = 0$, ce qui n'est pas. On peut aussi s'en assurer sur 6
\[
\boxed{\ \mathrm{sn}(-z) = -\mathrm{sn}\,z\ }
\]
\note{« 3 ou 4 », « 6 » : les formules (3), (4) et (6) ; dans la phrase, les
indices $\omega_1$, $\omega_2$ se lisent mal}

Sur 3 ou 4, on voit $\mathrm{sn}^2\omega_2 = \frac{1}{k}$ ;\note{le dénominateur est
écrit sur un autre signe} en choisissant convenablement le signe de $\omega_2$,
on aura
\[
\mathrm{sn}\,\omega_2 = \frac{1}{?}\,,
\]
\note{le dénominateur est surchargé et ne se lit pas ; la formule est
soulignée}
l'équation différentielle vérifiée par $\mathrm{sn}\,z$ donne
\[
\mathrm{sn}'\omega_2 = 0 \qquad \text{de même} \qquad \mathrm{sn}\,{-\omega_2} = -1
\quad \mathrm{sn}'\,{-\omega_2} = 0
\]
\note{le « $1$ » de $\mathrm{sn}\,{-\omega_2} = -1$ est une grosse tache d'encre}
On voit alors sur 4 que\struck{\ill{}}, puis sur l'équation différentielle, que
\[
\mathrm{sn}(\omega_1+\omega_2) = \frac{1}{k} \quad \text{soit} \quad
\mathrm{sn}\,\omega_3 = -\frac{1}{k} \qquad \mathrm{sn}'\omega_3 = 0 .
\]

\page{68}
On pose
\[
(8)\ \left\lbrace
\begin{aligned}
\mathrm{cn}^2 z &= 1 - \mathrm{sn}^2 z = \frac{\wp z - e_2}{\wp z - e_1}
= -\frac{1}{k^2}\bigl[\wp(z-\omega_1) - e_3\bigr]
&\qquad& \boxed{\mathrm{cn}\,0 = +1}\\
\mathrm{dn}^2 z &= 1 - k^2\mathrm{sn}^2 z = \frac{\wp z - e_3}{\wp z - e_1}
= -\bigl[\wp(z-\omega_1) - e_2\bigr]
&& \boxed{\mathrm{dn}\,0 = +1}
\end{aligned}
\right.
\]
\note{avant $-\frac{1}{k^2}[\ldots]$, un premier « $-\wp(z-\omega_1)$ » est biffé ; dans $\mathrm{cn}\,0 = +1$ le signe
est surchargé}
On a
\[
(9)\ \left\lbrace
\begin{aligned}
\mathrm{cn}\,z &= \frac{\sigma\omega_1}{\sigma\omega_2}\,e^{(\eta_2-\eta_1)z}\,
\frac{\sigma(z-\omega_2)}{\sigma(z-\omega_1)}
= \frac{\sigma\omega_1}{\sigma\omega_2}\,e^{-(\eta_2-\eta_1)z}\,
\frac{\sigma(z+\omega_2)}{\sigma(z+\omega_1)}\\
\mathrm{dn}\,z &= \frac{\sigma\omega_1}{\sigma\omega_3}\,e^{(\eta_3-\eta_1)z}\,
\frac{\sigma(z-\omega_3)}{\sigma(z-\omega_1)}
= \frac{\sigma\omega_1}{\sigma\omega_3}\,e^{-(\eta_3-\eta_1)z}\,
\frac{\sigma(z+\omega_3)}{\sigma(z+\omega_1)}
\end{aligned}
\right.
\]
\note{les indices des $\omega$ dans les quotients $\frac{\sigma\omega}{\sigma\omega}$
sont petits et se lisent mal}

\[
(10)\ \left\lbrace
\begin{array}{lll}
{}^{*}\ \mathrm{cn}(-z) = \mathrm{cn}\,z & \mathrm{dn}(-z) = \mathrm{dn}\,z &
\text{$\mathrm{cn}\,z$ et $\mathrm{dn}\,z$ sont paires}\\[4pt]
{}^{*}\ \mathrm{cn}(z+2\omega_1) = -\mathrm{cn}\,z & \mathrm{dn}(z+2\omega_3) = -\mathrm{dn}\,z & \\
\phantom{{}^{*}}\ \mathrm{cn}(z+2\omega_3) = -\mathrm{cn}\,z & \mathrm{dn}(z+2\omega_1) = -\mathrm{dn}\,z & \\[4pt]
{}^{*}\ \mathrm{cn}(z+2\omega_2) = \mathrm{cn}\,z & \mathrm{dn}(z+2\omega_2) = \mathrm{dn}\,z &
\text{$\mathrm{cn}$ admet pour périodes}\\
\phantom{{}^{*}}\ \mathrm{cn}(z+4\omega_1) = \mathrm{cn}\,z & \mathrm{dn}(z+4\omega_3) = \mathrm{dn}\,z &
\quad 2\omega_2,\ 4\omega_1,\ 4\omega_3\\
\phantom{{}^{*}}\ \mathrm{cn}(z+4\omega_3) = \mathrm{cn}\,z & \mathrm{dn}(z+4\omega_1) = \mathrm{dn}\,z &
\text{$\mathrm{dn}$ admet pour périodes}\\
& & \quad 2\omega_3,\ 4\omega_1,\ 4\omega_2
\end{array}
\right.
\]
\note{les astérisques sont de sa main, en tête de trois lignes ; plusieurs
indices des $\omega$ (dans les lignes en $\mathrm{dn}$ surtout) se lisent mal, et
dans la liste des périodes de $\mathrm{dn}$ les indices sont surchargés}
$\mathrm{cn}\,z$ et $\mathrm{dn}\,z$ sont fonctions elliptiques d'ordre $2$ ;

\emph{Pôles de $\mathrm{cn}\,z$}\note{souligné deux fois} homologues de $\omega_2$,
résidu $\frac{i}{k}$, et \struck{$\omega_2$} \add{$-\omega_2$} résidu
$-\frac{i}{k}$
\struck{si $\omega_1$, $\omega_2$ est direct}\add{\struck{(si $\omega_2$) \ill{}}}
\note{l'ajout au-dessus de la ligne est entouré et biffé ; « direct » est biffé
et souligné}
(l'ambiguïté de signe est levée par 9, joint à 6)

\struck{Pôles} \emph{Zéros de $\mathrm{cn}\,z$}\note{souligné deux fois} homologues
de $+\omega_2$ et $-\omega_2$ ; on a
\[
\mathrm{cn}'\omega_2 = -\frac{e^{\eta_2\omega_2}}{\sigma'\omega_2}
\qquad (9,\ \text{puis } 6)\,;
\]
à partir de la relation de définition de $\mathrm{cn}\,z$, on montrera
$\mathrm{cn}'^2\omega_2 = -\mathrm{sn}'^2\omega_2$
\[
\mathrm{cn}'\omega_2 = \pm\sqrt{1-k^2}
\]
\note{ainsi sur la page : « pôles » et « zéros » de $\mathrm{cn}$ tous deux
homologues de $\omega_2$ ; les exposants $\eta_2\omega_2$ se lisent mal}
Enfin
\[
\mathrm{cn}\,\omega_3 = \frac{i}{k}\,\frac{e^{\eta_2\omega_2}}{\sigma'\omega_2}
= -\frac{i}{k}\,\mathrm{cn}'\omega_2 = \pm\frac{i}{k}\sqrt{1-k^2}
\qquad \mathrm{cn}'\omega_3 = 0
\]

\page{69}
\emph{Pôles de $\mathrm{dn}\,z$}\note{souligné deux fois} homologues de
$\omega_2$, résidu $i$, et $-\omega_2$, résidu $-i$\note{« $\omega_2$, résidu
$i$ » est souligné ; l'indice du premier $\omega$ se lit mal}

\struck{Zéros}\note{un premier mot, en tête, est surchargé}
\emph{Zéros de $\mathrm{dn}\,z$}\note{souligné deux fois} homologues de $\omega_3$,
\emph{\uncertain{dérivée} $k\,\mathrm{cn}\,\omega_3$}
$= i\,\frac{e^{\eta_2\omega_2}}{\sigma'\omega_2} = -i\,\mathrm{cn}'\omega_2
= \pm i\sqrt{1-k^2}$ et de $-\omega_3$, (dérivée \uncertain{opposée} \ill{})
\note{« dérivée $k\,\mathrm{cn}\,\omega_3$ » est souligné ; la lecture de la
parenthèse finale est douteuse}

Enfin
\[
\mathrm{dn}\,\omega_2 = -\mathrm{cn}'\omega_2 = \frac{e^{\eta_2\omega_2}}{\sigma'\omega_2}
= \pm\sqrt{1-k^2} \qquad \mathrm{dn}'\omega_2 = 0
\]
\note{ainsi sur la page : « $\mathrm{dn}\,\omega_2$ » ; l'indice du premier
$\omega$ se lit mal}

\emph{Formules de dérivation}
\[
(11)\ \left\lbrace
\begin{aligned}
\mathrm{sn}'z &= \mathrm{cn}\,z\;\mathrm{dn}\,z\\
\mathrm{cn}'z &= -\mathrm{sn}\,z\;\mathrm{dn}\,z\\
\mathrm{dn}'z &= -k^2\,\mathrm{sn}\,z\;\mathrm{cn}\,z
\end{aligned}
\right.
\]
Formules liant $\mathrm{sn}\,z$, $\mathrm{cn}\,z$, $\mathrm{dn}\,z$ deux à deux
\[
(12)\quad \boxed{\ \mathrm{cn}^2 z + \mathrm{sn}^2 z = 1 \qquad
\mathrm{dn}^2 z + k^2\mathrm{sn}^2 z = 1 \qquad
\mathrm{dn}^2 z - k^2\mathrm{cn}^2 z = 1 - k^2\ }
\]
\note{devant $\mathrm{dn}^2z + k^2\mathrm{sn}^2z$, un signe est barré de
hachures ; le « $\mathrm{cn}$ » de la première formule de (11) est écrit sur
une autre lettre, et le numéro « (12) » se lit « (11) »}

$\mathrm{cn}\,z$ et $\mathrm{dn}\,z$ vérifient des équations analogues à celle de
$\mathrm{sn}\,z$
\[
(13)\ \left\lbrace
\begin{aligned}
\mathrm{cn}'^2 z &= (1-\mathrm{cn}^2 z)(1-k^2+k^2\mathrm{cn}^2 z)
= (1-k^2)(1-\mathrm{cn}^2 z)\Bigl(1 + \frac{k^2}{1-k^2}\mathrm{cn}^2 z\Bigr)\\
\mathrm{dn}'^2 z &= (1-\mathrm{dn}^2 z)\bigl(\mathrm{dn}^2 z - (1-k^2)\bigr)
= -(1-k^2)(1-\mathrm{dn}^2 z)\Bigl(1 - \frac{1}{1-k^2}\mathrm{dn}^2 z\Bigr)
\end{aligned}
\right.
\]
\note{le second facteur de la première ligne, « $1-k^2+k^2\mathrm{cn}^2z$ »,
est lu sous réserve}

Aussi $\mathrm{cn}\,z$ et $\mathrm{dn}\,z$ s'exprimeraient-ils directement au moyen
de la fonction $\mathrm{sn}\,z$ qui correspond respectivement à
$\rho = \frac{k}{\sqrt{1-k^2}}$ et $\rho = \frac{1}{\sqrt{1-k^2}}$, la solution de
l'équation
\[
\Bigl(\frac{dZ}{dz}\Bigr)^2 = (1-Z^2)(1-\rho^2 Z^2) \qquad
\]
\note{ainsi sur la page, en deux morceaux : « la solution de l'équation »
est écrit à gauche, sous « la fonction $\mathrm{sn}\,z$ » ; « autre » est
biffé, et la page s'arrête là}

\page{70}
\note{la page reprend la fonction $\mathrm{sn}$ sous le nom $\mathrm{sn}_1$, avec
une nouvelle numérotation des formules, (1) à (4)}
\[
(1)\quad \mathrm{sn}_1^2 z = \frac{1}{\wp z - e_1} \qquad (\mathrm{sn}_1'0 = +1)
\]
\[
(2)\quad \mathrm{sn}_1 z = \sigma\omega_1\,e^{\eta_1 z}\,\frac{\sigma z}{\sigma(z+\omega_1)}
= -\sigma\omega_1\,e^{-\eta_1 z}\,\frac{\sigma z}{\sigma(z-\omega_1)}
\]
\struck{$\mathrm{sn}_1^2 z = \frac{1}{\wp z - e_1}$}
\[
(3)\quad \mathrm{sn}_1 z\;\mathrm{sn}_1(z+\omega_1) = \frac{1}{k_1} \qquad
\mathrm{sn}_1^2 z = \frac{1}{k_1^2}\bigl(\wp(z-\omega_1) - e_1\bigr) \qquad
\begin{array}{l}
k_1 = -\frac{e^{\eta_1\omega_1}}{\sigma'\omega_1}\\[2pt]
k_1^2 = (e_3-e_1)(e_2-e_1)
\end{array}
\]
\note{l'indice de $\omega$ dans l'exposant de $k_1$ est surchargé ; le second
facteur de $k_1^2$ se lit « $(e_3-e_1)$ » une seconde fois, un indice étant
peu net}
\[
(4)\ \left\lbrace
\begin{array}{l}
\mathrm{sn}_1(z+2\omega_1) = \mathrm{sn}_1 z\\
\mathrm{sn}_1(z+2\omega_2) = -\mathrm{sn}_1 z\\
\mathrm{sn}_1(z+2\omega_3) = -\mathrm{sn}_1 z\\[4pt]
\mathrm{sn}_1(z+4\omega_2) = \mathrm{sn}_1 z\\
\mathrm{sn}_1(z+4\omega_3) = \mathrm{sn}_1 z\\[4pt]
\mathrm{sn}_1(-z) = -\mathrm{sn}_1 z\\[4pt]
\mathrm{sn}_1(2\omega_2 - z) = \mathrm{sn}_1 z\\
\mathrm{sn}_1(2\omega_3 - z) = \mathrm{sn}_1 z
\end{array}
\right.
\]
\note{le numéro (4) est souligné}
$\mathrm{sn}_1 z$ a les périodes $2\omega_1$, $4\omega_2$, $4\omega_3$. La
substitution $z\to z+2\omega_2$ et $z\to z+2\omega_3$ équivaut à un changement de
signe. $\mathrm{sn}_1 z$ est impair. Les substitutions $z\to 2\omega_2 - z$ et
$z\to 2\omega_3 - z$ laissent $\mathrm{sn}_1 z$ invariante.

$\mathrm{sn}_1 z$, fonction elliptique ayant les périodes $2\omega_1$, $4\omega_2$,
$4\omega_3$, d'ordre deux

\emph{pôles de $\mathrm{sn}_1 z$} les homologues de \emph{$\omega_1$, résidu
$+\frac{1}{k_1}$}, et $\omega_1+2\omega_2$ ou $\omega_1+2\omega_3$, résidu
$-\frac{1}{k_1}$

\emph{zéros de $\mathrm{sn}_1 z$} homologues de \emph{$0$, dérivée $+1$}, et
$2\omega_2$ ou $2\omega_3$, dér.\ $-1$
\[
\left\|\
\begin{array}{ll}
\mathrm{sn}_1\omega_2 = \frac{\sigma\omega_1\,\sigma\omega_2}{\sigma(\omega_1+\omega_2)}\,
e^{\eta_1\omega_2} = \pm\frac{1}{\sqrt{e_2-e_1}} & \mathrm{sn}_1'\omega_2 = 0\\[6pt]
\mathrm{sn}_1\omega_3 = \frac{\sigma\omega_1\,\sigma\omega_3}{\sigma(\omega_1+\omega_2)}\,
e^{\eta_1\omega_3} = \pm\frac{1}{\sqrt{e_3-e_1}} & \mathrm{sn}_1'\omega_3 = 0
\end{array}
\right.
\]
\note{les deux lignes sont marquées d'un double trait vertical à gauche ; la
dernière est coupée par le bas de la feuille, et son dénominateur se lit
$\sigma(\omega_1+\omega_2)$}

\page{71}
On pose
\[
(5)\ \left\lbrace
\begin{aligned}
\mathrm{cn}_1^2 z &= 1 - (e_2-e_1)\,\mathrm{sn}_1^2 z = \frac{\wp z - e_2}{\wp z - e_1}
&\qquad& (\mathrm{cn}_1 0 = 1)\\
\mathrm{dn}_1^2 z &= 1 - (e_3-e_1)\,\mathrm{sn}_1^2 z = \frac{\wp z - e_3}{\wp z - e_1}
&& (\mathrm{dn}_1 0 = 1)
\end{aligned}
\right.
\]
\note{le dernier dénominateur est lu $\wp z - e_1$ ; l'indice est peu net}
(6)\quad \emph{pôles} : les homologues de $\omega_1$, résidu $\frac{1}{k}$ ;
et $\omega_1+2\omega_2$, résidu $-\frac{1}{k}$

\emph{zéros} : les homologues de $0$, dérivée $+1$ ; et $2\omega_2$,
dérivée $-1$
\note{les deux lignes sont réunies par une accolade sous le numéro 6 ;
« pôles », « zéros », $\omega_1$ et $0$ sont soulignés deux fois}
\note{le $2$ de $2\omega_2$ est repassé à l'encre}
En choisissant convenablement $\omega_2$ :
\[
\underline{\ \mathrm{sn}\,\omega_2 = 1 \quad \mathrm{sn}'\omega_2 = 0 \qquad
\mathrm{sn}\,\omega_3 = -\frac{1}{k} \quad \mathrm{sn}'\omega_3 = 0\ }
\]

\page{73}
\[
(8)\ \left\lbrace
\begin{aligned}
\mathrm{sn}_1'z &= \mathrm{cn}_1 z\;\mathrm{dn}_1 z\\
\mathrm{cn}_1'z &= -(e_2-e_1)\,\mathrm{sn}_1 z\;\mathrm{dn}_1 z\\
\mathrm{dn}_1'z &= -(e_3-e_1)\,\mathrm{sn}_1 z\;\mathrm{cn}_1 z
\end{aligned}
\right.
\]
$\mathrm{sn}_1 z$, $\mathrm{cn}_1 z$, $\mathrm{dn}_1 z$ vérifient des équations
différentielles de type analogue. On a notamment
\[
(9)\quad \mathrm{sn}_1'^2 z = \bigl(1-(e_2-e_1)\,\mathrm{sn}_1^2 z\bigr)
\bigl(1-(e_3-e_1)\,\mathrm{sn}_1^2 z\bigr)
\]
\note{le reste de la page est blanc}

\section*{Fonctions elliptiques}
\note{titre encadré, de sa main, en tête de la page 74 ; sous lui, souligné
deux fois, le sous-titre « Expression des fonctions elliptiques par $\wp u$,
$\wp' u$ », relié par un trait au mot « généralités » de la ligne suivante}

\page{74}
\emph{1) généralités.} On appelle \emph{fonction elliptique} une fonction
méromorphe $f$ dans $\mathbf{C}$ qui admet deux périodes $\omega_1$, $\omega_2$
linéairement ind.\ sur $\mathbf{R}$ ; si $f$ n'est pas constante, le groupe des
périodes (sous-groupe fermé de $\mathbf{C}$ ($=\mathbf{R}^2$)) est discret, soit
$\mathfrak{P}$ \add{et de rang $2$}. On étudie \add{l'ensemble} des fonctions
elliptiques qui admettent $\mathfrak{P}$ comme groupe des périodes.
\note{la lettre notée $\mathfrak{P}$ est une grande majuscule gothique, qu'on
rend ainsi dans tout le lot ; l'ajout « et de rang $2$ » est écrit au-dessus de
la ligne, « l'ensemble » est appelé par un trait}
\struck{\ill{}}, soit $\lbrace 2\omega_1, 2\omega_2\rbrace$ une base de $\mathfrak{P}$
sur $\mathbf{Z}$ (définie à une substitution entière de déterminant $\pm 1$ près)
on étudie l'ens.\ des fonctions méromorphes telles que
\[
(1)\quad f(z+2\omega_1) = f(z+2\omega_2) = f(z)
\]
Il forme un \emph{corps} $E$, qu'on peut évidemment identifier au corps des
fonctions méromorphes sur la variété analytique quotient $\mathbf{C}/\mathfrak{P}$,
qui est topologiquement isomorphe au tore. -- On sait, d'après la théorie
\add{de Riemann,} \ill{}\note{un mot souligné en tête de ligne, suivi d'un
crochet fermant} que \add{pour} \struck{\ill{}} \struck{toute} $f\in E$ non
constante (et il en existe) $E$ est de degré fini sur $\mathbf{C}(f)$, i.e.\
\struck{\ill{}} toute fonction elliptique $g\in E$ est \struck{\ill{}} solution
d'une équation $P(f,g)=0$, $P$ polynôme en $f$ et $g$ \add{etc.}\ \ill{}
\ill{}. Nous obtiendrons directement des résultats plus précis.

\struck{Une fonction elliptique sans pôle est une constante}
\note{la phrase est rayée d'un trait épais}

\struck{a) Existence d'une f\ill{}}

a$'$) Somme des résidus nulle, \add{nb.\ de} zéros égal au nb.\ de pôles
\add{comptés avec leur multiplicité} ; \ldots\ -- immédiat par Liouville
\note{ajouts serrés en interligne, lus sous réserve ; « immédiat par
Liouville » est souligné}

b) Existence d'une fonction elliptique
\[
(2)\quad \wp(z) = \frac{1}{z^2} + \mathop{{\sum}'}_{m,n}\left(
\frac{1}{(z-2m\omega_1-2n\omega_2)^2} - \frac{1}{(2m\omega_1+2n\omega_2)^2}\right)
\]
définit une fonction holomorphe dans \struck{\ill{}} le plan, ayant
\struck{\ill{}} $\mathfrak{P}$ comme ensemble des pôles (pôles doubles).
On a
\[
(3)\quad \wp'(z) = -\frac{2}{z^3} - 2\mathop{{\sum}'}_{m,n}
\frac{1}{(z-2m\omega_1-2n\omega_2)^3} = -2\sum
\frac{1}{(z-2m\omega_1-2n\omega_2)^3}
\]

\page{75}
On voit directement que $\wp'(z)$ est elliptique \struck{\ill{}} changeant
$z$ en \struck{$z+2\omega_1$}.\note{le mot biffé, en fin de ligne, se lit mal}
De $\wp'(z) - \wp'(z+2\omega_i) = 0$ on déduit
$\wp(z) - \wp(z+2\omega_i) = \text{c}^{\text{te}}$, cette constante est
\uncertain{nulle} en particulier $z=-\omega_i$ : $\wp(-\omega_i) - \wp(\omega_i)$,
or $\wp$ est paire, donc la constante est nulle. \ill{} \ill{}.
\note{le raisonnement se lit sous réserve ; les indices $i$ sont peu nets}

Il existe une fonction elliptique et une seule (de périodes $2\omega_1$ et
$2\omega_2$) qui n'ait comme pôles \add{que les points de}
$2\mathfrak{P}$\note{ainsi sur la page, semble-t-il : « $2\mathfrak{P}$ »},
\struck{et \ill{} pôle \ill{} \ill{}} \add{et telle que} $\wp(z) - \frac{1}{z^2}$
soit holomorphe \add{et nulle} à l'origine. \struck{\ill{} \ill{} \ill{}
\ill{} \ill{} \ill{}} \ill{} la unicité $\wp(u)$, et de \ill{} découle par
l'expression de Mittag-Leffler (2), cette fonction est paire,
\struck{et \ill{} $\wp(z) - \frac{1}{z^2}$ est nulle} :
\[
\wp z = \frac{1}{z^2} + c_1 z^2 + c_2 z^4 + \cdots
\]
(L'unicité résulte immédiatement de a) ; la propriété \ill{} \ill{}
immédiat)
\note{tout ce paragraphe est marqué d'un trait vertical à gauche ; le passage
biffé l'est d'un trait ondulé}

c) \emph{Toute fonction elliptique} \add{paire $f$} \emph{\uncertain{n'ayant}
de pôles que aux points de $\mathfrak{P}$, est un polynôme en $\wp$ et un
seul}\note{ainsi sur la page, semble-t-il} -- \emph{Le polynôme est
unitaire}.\note{l'énoncé est souligné ligne à ligne}

La réciproque est évidente, ainsi que l'unicité, $\wp$ étant \uncertain{transcendant}
sur $\mathbf{C}$. -- $f$ \ill{} \ill{} \ill{} d'ordre pair $2k$ \ill{} \ill{}
pair $0$, \ill{} \ill{}, \ill{} \ill{} \ill{} \ill{} $k$ (évident si
$k=0$). \uncertain{Or}
\struck{$\frac{f(z)}{\wp(z)}$} \struck{$c_k$} si
\[
f(z) = \frac{c_k}{z^{2k}} + \cdots, \qquad f(z) - c_k\,\wp(z)^k
\]
sera
\emph{une fonction elliptique ayant \ill{} \ill{} \ill{} fonction elliptique
paire, n'ayant pour \ill{} \ill{} d'ordre $2(k-1)$ au point $0$}, \ill{} \ill{}
\ill{} \ill{} de pôles que dans $\mathfrak{P}$, et $0$ pour pôle d'ordre
$\leqslant$\note{lu sous réserve} $2(k-1)$, \ill{} \ill{} \ill{} \ill{} la
démonstration.
\emph{Notons que le degré de $P$, où $f = P(\wp)$, est $k$}
\note{plusieurs lignes
de ce passage sont soulignées}

d) \emph{Toute fonction elliptique paire $f$ \ill{} \ill{} fonction rationnelle
en $\wp$, et \uncertain{réciproquement}.} -- \emph{unicité}

Soit $\mathfrak{P}$\note{ici une lettre en forme de $\mathfrak{P}$ plus petite,
peut-être autre} le \uncertain{polynôme} \ill{} \ill{} \ill{} par $f$,
soient $a_1, \ldots, a_n$ \add{\uncertain{l'ensemble} des} \struck{\ill{} \ill{}
\ill{}} pôles de $f$ \add{\ill{} \ill{} $P$} \ill{} \ill{} $\mathfrak{P}$,
\struck{\ill{} \ill{} \ill{}}

\page{76}
\note{la page entière est barrée de longs traits obliques, et plusieurs lignes
sont en outre rayées d'un trait horizontal ; elle se lit par fragments, qu'on
donne ici sans marquer chaque rature}
\ldots\ de la forme $\wp' R_2(\wp)$, $R_2$ fonction rationnelle en $\wp$,
\ill{} \ill{} \ill{} que le dénominateur $Q$ de $R_2$ est une constante.
Autrement, $Q(\wp)$ \add{$=\prod(\wp - a_i)$} \uncertain{aurait} des zéros hors
de $\mathfrak{P}$, il faudrait que ces zéros fussent zéros de $\wp'$, avec une
multiplicité au moins égale ; or $\wp'$ a \ldots

\uncertain{Prouvons} d'abord le lemme :

\emph{\struck{la} Une fonction elliptique a au moins deux pôles non congruents
mod $\mathfrak{P}$, ou un pôle double.}
\struck{On peut \ill{} \ill{} \ill{} \ill{} \ill{} \ill{}}
\marginal{\ill{} \ill{} \ill{} \ill{} \ill{} \ill{} \ill{} \ill{}}
\note{deux lignes illisibles à gauche du paragraphe suivant, encadrées d'un
trait}
\ldots\ \ill{} de deux \emph{\uncertain{simples}} a un degré \ill{} \ill{}. Plus
\uncertain{élémentairement}, s'il n'y avait qu'un pôle \ill{} \ill{} congruent
mod $\mathfrak{P}$, \ill{} \ill{} \ill{} \ill{} \ill{} \ill{} pôle \ill{} de
$\mathfrak{P}$, \ill{} \ill{} \ill{} \ill{} $0$ \ill{} \ill{} pôle
\struck{$f = \frac{1}{z} + \cdots$ \ill{}}, \ill{} \ill{} \ill{} \ill{}
\ldots\ $\wp'$ serait fonction elliptique impaire, donc \ill{} \ill{}
$\wp^3 = \text{c}^{\text{te}}\,\wp'$ (b), \ldots\ \ill{} \ill{} \ill{} $f\in
\mathbf{C}(\wp)$, $f^2 = P$.

\ldots\ $\wp'$ aurait \struck{\ill{}} de zéros (dans $\mathbf{C}/\mathfrak{P}$)
\ldots\ \emph{\uncertain{jamais}} $3$, d'ailleurs \ill{} \ill{}, comme \ill{}
\ill{} \ill{}, \ldots\ $\omega_1$, $\omega_2$, $\omega_3$ sont des pôles \ldots

\marginal{\ill{} \ill{} des \ill{} \ill{} \ill{} ; \ill{} \ill{} \ill{}
$=\mathbf{C}/\mathfrak{P}$ \ill{} \ill{} \ill{} deux fois \ill{} \ill{} \ill{}
\ill{}}
\note{note en biais dans la marge gauche, à l'encre, bordée d'un trait
vertical, lue par fragments}
$(2)$ que $\wp'$ fonction elliptique impaire n'ayant pas de pôles en
$\omega_1$, $\omega_2$, $\omega_3 = \omega_1+\omega_2$) $\wp'$ s'annule en ces
points ($f(\omega_i) = -f(-\omega_i)$, or $f(-\omega_i) = f(\omega_i)$ \ldots\
période). Il suit, \ldots\ \emph{que $\wp'$ a $\omega_1$, $\omega_2$, $\omega_3$
comme zéros simples} (\ldots\ \ill{} dans $\mathfrak{P}$), \ldots\ $\wp'$ \ldots\
\struck{\ill{}} de $2$ \ldots\ $f^3$.

\page{77}
\struck{\ill{} $2\omega_1$ \ill{} $2\omega_2$ \ill{} \ill{} \ill{}}
\note{la première ligne, qui achève le passage barré de la page 76, est rayée}
$l_i$ la multiplicité \ill{} $a_i$. Alors
\[
g(z) = f(z)\prod_{1\leqslant i\leqslant n}\bigl[\wp(z) - \wp(a_i)\bigr]^{l_i}
\]
est fonction elliptique paire, n'ayant pour pôles que des points de
$\mathfrak{P}$. Il suit donc d'ailleurs c).
\struck{\ill{} \ill{} \ill{} \ill{} $f$ \ill{} \ill{} \ill{} \ill{} \ill{}
fonction elliptique paire}
\note{deux lignes rayées et biffées d'un trait ondulé ; une première tentative
de « e) », également biffée, les précède}

e) \emph{Le corps $E$ des fonctions elliptiques est de degré $2$ sur
$\mathbf{C}\langle\wp\rangle$.} En effet, $\mathbf{C}\langle\wp\rangle$ est le
corps des invariants de l'automorphisme $f\to\check{f}$ de $E$, qui est d'ordre
$2$ (notons que $E\neq\mathbf{C}\langle\wp\rangle$, car $\wp'$ est
\emph{impaire}) \uncertain{comme} \uncertain{nous} \uncertain{verrons}
$\wp'\notin\mathbf{C}\langle\wp\rangle$.
\note{$\check{f}$ désigne manifestement $z\mapsto f(-z)$ ; le signe $\notin$ est
lu d'après un $\in$ surchargé}

f) \emph{Toute fonction elliptique $f$ se met d'une façon et d'une seule sous la
forme}
\[
(3)\quad f = R_1(\wp) + \wp' R_2(\wp)
\]
\emph{où $R_1$ et $R_2$ sont des fonctions rationnelles.}

\struck{\ill{}} N.B., \uncertain{comme}

g) \struck{Le nombre de pôles} \add{Pour que la fonction $f$ de (3) ait ses
pôles} \emph{\uncertain{contenus} dans $\mathfrak{P}$, il faut et il suffit que
$R_1$ et $R_2$ soient des polynômes} ; \emph{alors on a au moins pôle
double}.\note{ainsi sur la page, semble-t-il ; l'énoncé est souligné}

\struck{Il suffit de prouver que si le pôle de $f$ \ldots\ \ill{} dans
$\mathfrak{P}$, $f$ est de la forme $P_1(\wp) + \wp' P_2(\wp)$, $P_1$ et $P_2$
polynômes. $\mathcal{E}$ \uncertain{écrivons} $f = \frac{f+\check{f}}{2} +
\frac{f-\check{f}}{2}$, on est ramené au cas où $f$ est paire -- traité dans
c) -- ou au cas où $f$ est impaire. Mais sur (3), on voit que \ill{}
\ill{} \ill{} que $f$ \ldots}
\note{ce dernier paragraphe est encadré et barré de trois traits obliques}

\page{78}
La \uncertain{première} de g) est \struck{\ill{}} \add{\uncertain{multiplicité}}
\note{ajout en interligne, entre crochets} facile (réciproquement, \ldots\
\struck{\ill{}} du pôle à l'origine) puisque la suite est au moins d'ordre $2$)
\ldots\ \ill{} \add{aussi} facilement que \struck{\ill{}} fonction elliptique
impaire est de la forme $\wp' R(\wp)$ -- $R$ fonction rationnelle -- (évidemment
d'après b)) et où $R$ est un polynôme si et seulement si l'ens.\ des pôles est
$\subset\mathfrak{P}$ (\uncertain{c.-à-d.}\ \ill{} g)).
\note{paragraphe écrit vite ; plusieurs mots lus sous réserve}

\emph{2) Les relations entre $\wp$, $\wp'$}
\note{titre encadré et souligné deux fois ; le numéro « 2) » est repassé sur
un autre chiffre, entouré}
\[
(1)\ \left\lbrace
\begin{aligned}
\wp &= \frac{1}{z^2} + \mathop{{\sum}'}_{m,n}\left(
\frac{1}{(z-2m\omega_1-2n\omega_2)^2} - \frac{1}{(2m\omega_1+2n\omega_2)^2}\right)\\
\wp' &= -2\sum_{m,n}\frac{1}{(z-2m\omega_1-2n\omega_2)^3}
\end{aligned}
\right.
\]
\note{après la seconde ligne, « $= -\frac{2}{z^3} + \cdots$ » biffé de
hachures}
\[
(2)\quad \wp z = \frac{1}{z^2} + c_1 z^2 + c_2 z^4 + \cdots
\]
Les $c_i$ s'obtiennent \uncertain{aisément} \add{à partir de} (1), en
dérivant les deux membres, \ill{} \ill{} \ill{} \ill{} \ill{} de $\wp z$,
\ill{} \ill{} \ill{} pour $z\to 0$
\[
(3)\ \left\lbrace
\begin{array}{l}
c_1 = 3\displaystyle\mathop{{\sum}'}_{m,n}\frac{1}{(2m\omega_1+2n\omega_2)^4}\\
\quad -\ -\ -\\
c_k = (2k+1)\displaystyle\mathop{{\sum}'}_{m,n}\frac{1}{(2m\omega_1+2n\omega_2)^{2k+2}}
\end{array}
\right.
\]
\note{le premier terme des dénominateurs de $c_1$ est surchargé}

Les zéros de $\wp'(z)$ dans $\mathfrak{P}$ sont $\omega_1$, $\omega_2$,
$\omega_3 = \omega_1+\omega_2$. On pose
\[
\boxed{\ \wp\omega_i = e_i\ }
\]
ce sont racines doubles de l'équation $\wp z - e_i = 0$ ; il suit que les $e_i$
sont distincts, car $\wp z - e_i$ n'a que $2$ racines (distinctes ou non) dans
$\mathbf{C}/\mathfrak{P}$.
\note{ainsi sur la page : « dans $\mathfrak{P}$ » ; $e_i$ et « sont distincts »
sont soulignés}

\page{79}
\emph{Équations différentielles de $\wp u$.} Il \uncertain{faut} remarquer, que
\ill{} \ill{} deux fonctions \ill{} elliptiques, $f$ et $f'$ \ill{} \ill{} liées
par une relation polynomiale -- \struck{\ill{}} \emph{la fonction elliptique
\uncertain{satisfait} à une équation différentielle polynomiale du premier
ordre}. -- Pour $\wp$, on voit que $\wp'^2$ est de degré $2$ dans
$\mathbf{C}\langle\wp\rangle$, \ill{} ; \ill{} pour \ill{} \ill{} relation
\uncertain{considérée} : $\wp'^2$ est \emph{fonction} elliptique paire, de degré
$6$, donc \ldots\ $\wp$ de degré $3$, \ldots\ \uncertain{soit}
\add{\uncertain{par identification}} \struck{\ill{}} \ldots\ polynôme est de
\struck{\ill{}}
\[
\wp'^2 = 4\wp^3 - 20c_1\,\wp - 28c_2
\]
\note{un signe devant $20c_1$ est effacé d'une tache d'encre ; la dernière
lettre, après $28c_2$, est un $\wp$ barré}
soit
\[
(?)\quad \boxed{\ \wp'^2 = 4\wp^3 - g_2\,\wp - g_3\ } \qquad
(g_2 = 20c_1,\ g_3 = 28c_2)
\]
\note{le numéro, petit, se lit mal (un $4$ ou un $7$) ; celui de la formule
encadrée suivante (un $5$ ou un $8$) aussi ; la page 80 poursuit par (6), ce
qui ferait attendre (4) et (5). Les deux derniers $\wp$ du cadre sont surchargés}
\struck{Le calcul (\ldots\ \ill{} \ill{} \ill{} multiplicité de $\wp'$)
\ill{} \ill{} \ill{} \ill{}}, on \uncertain{remarquera} que
\[
\frac{\wp'^2}{(\wp-e_1)(\wp-e_2)(\wp-e_3)}
\]
est fonction elliptique holomorphe, (car \ill{} \ill{} \ill{} \ill{} \ill{}
\ill{} \ill{} à l'origine, \ill{} \ill{} \ill{} \ill{} $e_i$) \ill{}, donc
\ill{} \ill{} \ill{} ; \uncertain{elle} est égale à $4$ (\uncertain{regarder}
les développements de $\wp$ et $\wp'$ à l'origine) d'où
\[
(?)\quad \boxed{\ \wp'^2 = 4(\wp-e_1)(\wp-e_2)(\wp-e_3)\ }
\]
\ill{} \ill{} \ill{} \ill{} par la \uncertain{polynôme} -- \ill{}
\uncertain{précédent} est $4(\wp^3-e_1)(\wp-e_2)(\wp-e_3)$.\note{ainsi sur la
page : $\wp^3$ au premier facteur}
\struck{$\Sigma e_i = 0$, \ill{} \ill{} \ill{} \ill{} \ill{} \ill{} $\wp'$
\ill{} \ill{}}
\note{la dernière ligne est biffée}

\page{80}
En poussant plus loin le calcul d'identification, on obtiendrait les
\uncertain{expressions} de $c_3$, $c_4$ \ldots\ en fonction de $c_1$, $c_2$. Le
calcul est d'ailleurs \uncertain{simplifié} \ill{} \ill{} \ill{} \ill{} \ill{}
\ill{} \ill{} \uncertain{dérivées} de (4) \uncertain{et} utilisons les
\uncertain{formules} \ill{}, par dérivation\note{deux mots barrés dans la
quatrième ligne}
\[
(6)\quad \boxed{\ \wp'' = 6\wp^2 - \tfrac{1}{2}g_2 \qquad \wp''' = 12\,\wp\wp'\ }
\]
Les relations \uncertain{précédentes} \ill{}
\[
\boxed{\ c_3 = \frac{c_1^2}{3}\,, \qquad c_4 = \frac{3}{11}\,c_1 c_2\ } \qquad (7)
\]

\emph{3) Fonction $\zeta(u)$}
\note{titre encadré ; le numéro « 3) » est écrit à la suite d'un chiffre
noirci}
\[
(7)\quad \zeta u - \frac{1}{u} = \int_0^u -\Bigl(\wp z - \frac{1}{z^2}\Bigr)dz
\qquad \zeta' u = -\wp u
\]
\[
(8)\quad \zeta u = \frac{1}{u} - \frac{c_2}{3}u^3 - \cdots - \frac{c_n}{2n+1}u^{2n+1} - \cdots
\]
\note{ainsi sur la page : le premier coefficient se lit $\frac{c_2}{3}$, avec un
indice surchargé ; on attendrait $\frac{c_1}{3}$}
$\zeta$ est impaire, admet les points de $\mathfrak{P}$ pour pôles simples, avec
pour résidu $1$.
\[
(9)\quad \zeta z = \frac{1}{z} + \sum_{\omega\in\mathfrak{P}^*}
\Bigl(\frac{1}{z-\omega} + \frac{1}{\omega} + \frac{z}{\omega^2}\Bigr)
\]
\note{le premier dénominateur de la parenthèse est noirci et récrit ; les $z$
de la formule sont surchargés sur d'autres lettres}
$\zeta'$ \uncertain{étant} \uncertain{une} \uncertain{elliptique} (\ill{}
\ill{} \ill{}), on \ill{}
\[
(10)\quad \boxed{\ \zeta(z+2\omega) - \zeta(z) = 2\eta\ } \]
($= \text{c}^{\text{te}}$ ; \ill{} dérivée de la première nulle)
\note{le numéro (10) est entouré}
\uncertain{en} \uncertain{faisant} \uncertain{à} \uncertain{nouveau}
\ill{} $z = -\omega$, on obtient
\[
(11)\quad \eta = \zeta(\omega)\,;
\]
\note{le numéro (11) est entouré}
$\eta_1$, $\eta_2$, $\omega_1$, $\omega_2$ sont liés par une relation ;
si $\omega_1$, $\omega_2$ sont \uncertain{en} \uncertain{disposition}
\uncertain{directe}, on a
\[
\eta_1\omega_2 - \eta_2\omega_1 = \frac{i\pi}{2}
\]
\note{la page s'arrête sur cette formule, au bord inférieur}
\end{document}
